\subsection{乘法群 R*}

由第三章 § 6 第 7 目可知，实直线的拓扑在 \(\mathbf{R}^*\) 上诱导的拓扑与 \(\mathbf{R}^*\) 的乘法群结构相容；由于 \(\mathbf{R}^*\) 是局部紧空间 \(\mathbf{R}\) 的开子集，由此推出 \(\mathbf{R}^*\) 是局部紧拓扑群（第一章 § 9 第 7 目，命题 13），因而是完备的（第三章 § 3 第 3 目，命题 4 的推论 1；这也可由第三章 § 6 第 8 目的命题 8 得出）；当然，这后一性质涉及的是 \(\mathbf{R}^*\) 上的乘法一致结构，而不是 \(\mathbf{R}^*\) 上由 \(\mathbf{R}\) 的加法一致结构诱导的一致结构。

函数 \(xy\) 把集 \(\mathbf{Q}_{+} \times \mathbf{Q}_{+}\) 映入 \(\mathbf{Q}_{+}\) ，因此它把 \(\mathbf{R}_{+} \times \mathbf{R}_{+}\) 映入 \(\mathbf{R}_{+}\) （第一章 § 2 第 1 目，定理 1）；换句话说，两个 \(\geqslant 0\) 的实数之积 \(\geqslant 0\) 。于是公式 \((-x)y = -xy\) 与 \((-x)(-y) = xy\) 表明，一个 \(\geqslant 0\) 的数与一个 \(\leqslant 0\) 的数之积 \(\leqslant 0\) ，而两个 \(\leqslant 0\) 的数之积 \(\geqslant 0\) ；由此推出

\[
| x y | = | x |. | y |\tag{1}
\]

（它也可以由 \(\mathbf{Q} \times \mathbf{Q}\) 上相应关系的延拓得到）。

若 \(x > 0\) 且 \(y > 0\) ，我们有 \(xy \neq 0\) ，因而 \(xy > 0\) ；同样，若 \(x < 0\) 且 \(y > 0\) ，则 \(xy < 0\) ；若 \(x < 0\) 且 \(y < 0\) ，则 \(xy > 0\) 。特别地，若 \(x \neq 0\) ，我们有 \(x^2 > 0\) ，因此实数的平方和除非每个数都为零，否则不可能为零。

若 \(x > 0\) 且 \(y \leqslant z\) （相应地 \(y < z\) ），我们有 \(xy \leqslant xz\) （相应地 \(xy < xz\) ）；换句话说，比率 \(> 0\) 的位似保持 \(\mathbf{R}\) 上的次序。由于 \((-x)y = -xy\) ，比率 \(< 0\) 的位似把 \(\mathbf{R}\) 上的序变为反序。

若 \(x > 0\) ，我们有 \(\mathbf{1} / x > 0\) ，因为 \(x.(\mathbf{1} / x) = \mathbf{1} > 0\) 。若 \(0 < x < y\) ，我们有 \(xy > 0\) ，于是 \(x.(\mathbf{1} / xy) < y.(\mathbf{1} / xy)\) ，即 \(\mathbf{1} / y < \mathbf{1} / x\) 。因此，映射 \(x \to \mathbf{1} / x\) 把集 \(\mathbb{R}_+^*\) （即 \(> 0\) 的实数组成的集）映到其自身上，它是严格递减的。

我们同样看出，函数 \(\mathbf{1} / x\) 在 \(] \leftarrow, o[, \text{ 且因此函数 } \frac{\mathbf{1}}{x - a}\) 在区间 \(] \leftarrow, a[ \text{ 与 } ]a, \rightarrow [.\) 的每一个内严格递减。

由上文推出，\(R_{+}^{*}\) 是乘法群 \(R^{*}\) 的一个子群；而且，序关系 \(x \leqslant y\) 与 \(R_{+}^{*}\) 的乘法群结构相容；换句话说，\(R_{+}^{*}\) 是一个全序群。

两个 \(\geqslant0\) 的实数之积 \(\geqslant0\) 这一事实可以表述为：R 是一个有序域；上述所有性质对一切有序域都成立。

命题 2. 乘法群 \(\mathbf{R}^{*}\) （ \(\neq 0\) 的实数组成的）是一个拓扑群，同构于它的子群 \(\mathbf{R}_{+}^{*}\) 与 \(\mathbf{U}_{0}\) 的积，其中

\[
\mathrm{U} _ {0} = \{- 1, + 1 \}.
\]

对每个 \(x \neq 0\) ，令 \(\operatorname{sgn} x\) 表示 \(\frac{x}{|x|}\) （ \(x\) 的符号）。函数 \(\operatorname{sgn}\) 是 \(\mathbf{R}^*\) 到 \(U_0\) 上的一个同态。我们有 \(x = |x| \operatorname{sgn} x\) ，且 \(x\) 的这个表成 \(\mathbf{R}_+^*\) 的一个元素与 \(U_0\) 的一个元素之积的分解是唯一的；因而 \(\mathbf{R}^*\) 的群结构是 \(\mathbf{R}_+^*\) 与 \(U_0\) 的群结构之积。另一方面，映射 \(x \to |x|\) 连续，且 \(x \to \operatorname{sgn} x = \frac{x}{|x|}\) 也连续，因为 \(x \neq 0\) 。由此即得结论。

我们通过令 sgn o = o 把函数 sgn 延拓到整个 R。

我们将在第五章（§ 4 第 1 目，定理 1）看到，拓扑群 \(\mathbf{R}_{+}^{*}\) 同构于加法群 \(\mathbf{R}\) ；这将完成对拓扑群 \(\mathbf{R}^{*}\) 的结构的确定。

